\documentclass{article}
\usepackage{geometry}
\geometry{a4paper,scale=0.8}
\usepackage[utf8]{inputenc}

\usepackage{amsmath, amssymb}
\usepackage{amsthm}
\usepackage{enumitem}
\usepackage{blkarray}
\newtheorem{remark}{Remark}
\newtheoremstyle{breakstyle}
    {5pt}
    {10pt}
    {\normalfont}
    {0pt}
    {\bfseries}
    {\newline\indent}
    {0pt}
    {}
\theoremstyle{breakstyle}
\newtheorem{exercise}{Exercise}
\DeclareMathOperator{\rad}{Rad}
\title{Chapter 1 - Rings and Ideals}
\author{Flandre}
\date{23 Jan 2026}

\begin{document}
\maketitle

\begin{exercise}
    Suppose $x^{r}=0$.
    Notice
    \begin{align*}
        (1+x)\left( 1-x+x^2-x^3+\ldots +(-1)^{r-1}x^{r-1} \right) 
        &= 1-x^{r} \\
        &= 1 \\
    \end{align*}
    So $1+x$ is a unit in $A$.
    \qed
\end{exercise}

\begin{exercise}
    \leavevmode \vspace{-\baselineskip}
    \begin{enumerate}[label=({\roman*})]
        \item 
            \begin{enumerate}[label=({\alph*})]
                \item $(\Rightarrow)$:
                    Let $g=b_0+b_1x+\ldots +b_{m}x^{m}$ be the inverse of $f$.
                    Now
                    $$
                    (a_0+a_1x+\ldots +a_{n}x^{n})(b_0+b_1x+\ldots +b_{m}x^{m})=1
                    $$
                    By checking the coefficients of $x^{n+m},x^{n+m-1},x^{n+m-2}\ldots $, we get
                    $$
                    \begin{cases}
                        a_{n}b_{m}=0\\
                        a_{n}^2b_{m-1}=0\\
                        a_{n}^3b_{m-2}=0\\
                        \vdots\\
                        a_{n}^{m}b_1=0\\
                        a_{n}^{m+1}b_0=0
                    \end{cases}
                    $$
                    On the other hand, by checking the constant term directly,
                     $$
                    a_0b_0=1
                    $$
                    Therefore $a_0$ is unit, and $a_{n}^{m+1}=a_{n}^{m+1}b_0a_0=0$.
                    So $a_{n}$ is nilpotent.

                    Moreover, $-a_{n}x^{n}$ is a nilpotent element in $A[x]$.
                    By \emph{Ex.1},
                    $$
                    f-a_{n}x^{n}=a_0+a_1x+\ldots +a_{n-1}x^{n-1}
                    $$
                    is a unit in $A[x]$, so by iterating the process above we'll see $a_{n-1},\ldots ,a_1$ are nilpotent.
                \item $(\Leftarrow)$: By \emph{Ex.1}, the sum of a unit and a nilpotent element is a unit.
                    Notice $a_1x,a_2x^2,\ldots ,a_{n}s^{n}$ are nilpotent in $A[x]$ while $a_0$ is a unit in $A[x]$,
                    so their sum $f$ is a unit in $A[x]$.
            \end{enumerate}
            \qed
        \item 
            \begin{enumerate}[label=({\alph*})]
                \item $(\Rightarrow):$
                    Consider the constant term of power of $f$, from which we know $a_0$ is nilpotnet.
                    Nilpotent elements are closed under addition, so
                    $$
                    a_1x+a_2x^2+\ldots +a_{n}x^{n}=f-a_0
                    $$
                    is nilpotent, i.e.
                    $$
                    a_1+a_2x+\ldots +a_{n}x^{n-1}
                    $$
                    is nilpotent.
                    By induction we can have $a_0,a_1,\ldots ,a_{n}$ being nilpotent.
                \item $(\Leftarrow):$
                    When $a_0,a_1,\ldots ,a_{n}$ are all nilpotent, let $k\in \mathbb{Z}^{+}$ be sufficiently large,
                    and obviously we can finally reach $f^{k}=0$.
            \end{enumerate}
            \qed
        \item 
            \begin{enumerate}[label=({\alph*})]
                \item $(\Rightarrow):$ 
                    Suppose we have $g=b_0+b_1x+\ldots +b_{m}x^{m}$ such that $fg=0$.

                    If $b_{m}=0$, we're done.
                    Otherwise, simmilar to \emph{(i)}, by checking the coefficients of $x^{m+n},x^{m+n-1},\ldots $ respectively,
                    we'll obtain
                    $$
                    \begin{cases}
                        a_{n}b_{m}=0\\
                        a_{n-1}b_{m}=0\\
                        a_{n-2}b_{m}^2=0\\
                        \vdots\\
                        a_{1}b_{m}^{n-1}=0\\
                        a_0b_{m}^{n}=0
                    \end{cases}
                    $$
                    So $f\cdot b_{m}^{n}=0$.
                \item $(\Leftarrow):$
                    Trivial.
            \end{enumerate}
        \item 
            \begin{enumerate}[label=({\alph*})]
                \item $(\Leftarrow):$
                    Let $g=b_0+b_1x+..+b_{m}x^{m}$.
                    Since $g$ is primitive, $\exists \lambda_0,\lambda_1,\ldots ,\lambda_{m}\in A,s.t.$ 
                    $$
                    \lambda_0b_0+\lambda_1b_1+\ldots +\lambda_{m}b_{m}=1
                    $$
                    Therefore
                    $$
                    \lambda_0b_0f+\lambda_1b_1f+\ldots +\lambda_{m}b_{m}f=f
                    $$
                    Also, suppose
                    $$
                    \mu_0a_0+\mu_1a_1+\ldots +\mu_{n}a_{n}=1
                    $$
                    Since
                    $$
                    \begin{array}{r@{\:}l@{\:}l@{\:}l@{\:}l}
                    fg = \lambda_0 \mu_0 b_0 a_0 & + \lambda_1 \mu_0 b_1 a_0 x & + \lambda_2 \mu_0 b_2 a_0 x^2 & + \dots & + \lambda_m \mu_n b_m a_n x^{m+n} \\
                                                 & + \lambda_0 \mu_1 b_0 a_1 x & + \lambda_1 \mu_1 b_1 a_1 x^2 & + \dots \\
                                                 &                             & + \lambda_0 \mu_2 b_0 a_2 x^2 & + \dots \\
                                                 &                             & \vdots
                    \end{array}
                    $$
                    So
                    $$
                    (\lambda_0\mu_0b_0a_0)+(\lambda_0\mu_1b_0a_1+\lambda_1\mu_0b_1a_0)+(\lambda_0\mu_2b_0a_2+\lambda_1\mu_1b_1a_1+\lambda_2\mu_0b_2a_0)+\ldots +(\lambda_{m}\mu_{n}b_{m}a_{n})=1
                    $$
                \item $(\Rightarrow):$ When $fg$ is primitive.
                    We prove by contradiction and suppose $g$ is not primitive.
                    Since $fg$ is primitive, we may assume
                    $$
                    b_0a_0\xi_0+(b_1a_0+b_0a_1)\xi_1+(b_2a_0+b_1a_1+b_0a_2)\xi_2+\ldots +b_{m}a_{n}\xi_{m+n}=1
                    $$
                    Therefore
                    \begin{align*}
                        &(a_0\xi_0+a_1\xi_1+\ldots +a_{n}\xi_{n})b_0\\
                        +&(a_0\xi_1+a_1\xi_2+\ldots +a_{n}\xi_{n+1})b_{1}\\
                        +&(a_0\xi_2+a_1\xi_3+\ldots +a_{n}\xi_{n+2})b_{2}\\
                        &\vdots\\
                        +&(a_0\xi_{m}+a_1\xi_{m+1}+\ldots +a_{n}+\xi_{m+n})b_{m}=1
                    \end{align*}
                    Contradictory.
            \end{enumerate}
        \qed
    \end{enumerate}
\end{exercise}

\begin{remark}
Consider the sum of terms in each of $n+m-1$ anti-diagonals
$$
\begin{blockarray}{cccccc}
 & \lambda_0 & \lambda_1 & \lambda_2 & \cdots & \lambda_m \\
\begin{block}{c[ccccc]}
  \mu_0 & b_0 a_0 & b_1 a_0 x & b_2 a_0 x^2 & \cdots & b_m a_0 x^m \\
  \mu_1 & b_0 a_1 x & b_1 a_1 x^2 & b_2 a_1 x^3 & \cdots & b_m a_1 x^{m+1} \\
  \mu_2 & b_0 a_2 x^2 & b_1 a_2 x^3 & b_2 a_2 x^4 & \cdots & b_m a_2 x^{m+2} \\
  \vdots & \vdots & \vdots & \vdots & \ddots & \vdots \\
  \mu_n & b_0 a_n x^n & b_1 a_n x^{n+1} & b_2 a_n x^{n+2} & \cdots & b_m a_n x^{m+n} \\
\end{block}
\end{blockarray}
$$
\end{remark}

\begin{exercise}
    We can prove it by induction. View $A[x_1,\ldots ,x_{n}]$ as a ring over $A[x_1,\ldots ,x_{n-1}]$ in an indeterminate $x_{n}$.
    \qed
\end{exercise}

\begin{exercise}
    Recall that the Jacobson radical is the intersection of all maximal ideals, while the nilradical is the intersection of all prime ideals.
    We'll prove that every prime ideal $p\subseteq A[x]$ is maximal:

    $$
    \forall \text{prime ideal }p,\text{ it's maximal }
    $$
    $\Leftrightarrow$ 
    $$
    \forall\text{ prime ideal }p,\text{ Krull dimension of }p\text{ is }0
    $$
    $\Leftrightarrow$
    $$
    \dim{A[x]}=0
    $$
    $\Leftrightarrow$(By \emph{[Algebraic Geometry - Hartshorne]Proposition 1.7}) 
    $$
    \dim{\left\{ \emptyset \right\} }=0\text{, }\left\{ \emptyset \right\} \text{in affine }1\text{-space }A
    $$
    Which is trivial.
\end{exercise}

\begin{remark}
    By induction, we'll get Jacobson radical being equal to nilradical in $A[x_1,x_2,\ldots ,x_{n}],\forall n \in \mathbb{Z}^{+}$
\end{remark}

\begin{exercise}
    \leavevmode \vspace{-\baselineskip}
    \begin{enumerate}[label=({\roman*})]
        \item 
            \begin{enumerate}[label=({\alph*})]
                \item $(\Rightarrow):$ Suppose $fg=1$, where
                    $$
                    g=\sum_{n=0}^{\infty} b_{n}x^{n} 
                    $$
                    By checking the constant term of $fg$, we get
                    $$
                    a_0b_0=1
                    $$
                    Thus $a_0$ is a unit in $A$.
                \item $(\Leftarrow):$ Notice
                    $$
                    (a_0+a_1x+\ldots )(a_0^{-1}-a_1a_0^{-2}x+(a_1^2a_0^{-3}-a_2a_0^{-2})+\ldots )=1
                    $$
                    Such construction of $f^{-1}$ can goes forever, so $f$ is a unit in $A[[x]]$.
            \end{enumerate}
        \item 
            \begin{enumerate}[label=({\alph*})]
                \item When $f$ is nilpotent, obviously $a_0$ is nilpotent.
                    Nilpotent elements closed under addition, so $f-a_0$ is nilpotnet, indicating $a_1$ being nilpotent.
                    Iterate the process above and we get $a_0,a_1,a_2,\ldots $ being nilpotent.
                \item No. Counterexample can be reached on \emph{MathStackExchange}.
            \end{enumerate}
        \item 
            \begin{enumerate}[label=({\alph*})]
                \item $(\Leftarrow):$
                    For $a_0$ in maximal ideal $I$, we can easily check that
                    $$
                    \mathfrak{I}_{I}=\left\{ b_0+b_1x+b_2x^2+\ldots \mid b_0\in I,b_1,b_2,\ldots \in A \right\} 
                    $$
                    Is a maximal ideal in $A[[x]]$.
                    Therefore $f=\sum_{n=0}^{\infty} a_{n}x^{n}\in \mathfrak{I}_{I} $.
                    
                    So when $a_0$ is in the Jacobson Radical of $A$, $f$ must be in the Jacobson Radical of $A[[x]]$.
                 \item $(\Rightarrow):$
                     We only need to prove the maximal ideals of $A[[x]]$ are all in the form
                     $$
                    \mathfrak{I}_{I}=\left\{ b_0+b_1x+b_2x^2+\ldots \mid b_0\in I,b_1,b_2,\ldots \in A \right\} 
                     $$
                     where $I$ is a maximal ideal of $A$.

                     Indeed, it's easy to see that the constant terms of all the elements in a arbitary ideal $\mathfrak{I}$ of $A[[x]]$ forms an ideal in $A$,
                     we denote then by $I$ and let $\mathfrak{I}_{I}=\mathfrak{I}$.
                     When $\mathfrak{I}_{I}$ is maximal, $I$ must be maximal in $A$, otherwise $\mathfrak{I}_{I}$ can expand to a larger ideal $\mathfrak{I}_{\overline{I}}$ where $\overline{I}$ is a maximal ideal containing $I$.

                     Hence the ideal of $A[[x]]$ is maximal if and only if the constant term ideal is maximal in $A$.
                     And $f$ is in a arbitary maximal ideal will indicate $a_0$ belongs to the correspond maximal ideal.
                     Deduction of Jacobson Radical follows immediately.
            \end{enumerate}
        \item Proved in $(iii)$ above.
        \item 
            \textbf{(UNFINISHED)}
            \emph{(I'm not sure whether I understood this problem correctly.)}

            Consider
            \begin{align*}
                \varphi:A&\longrightarrow A[[x]]\\
                a&\longrightarrow a+0x+0x^2+\ldots 
            \end{align*}
            And define $\varphi^{-1}$ by taking the constant term in the formal power series.
            For prime ideal $\mathfrak{p}\subseteq A$,
            $$
            \mathfrak{p}A[[x]]=\left\{ a_0+a_1x+a_2x^2+\ldots \mid a_0,a_1,\ldots \in \mathfrak{p} \right\} 
            $$
            We can prove it's a 
            % So its preimage is 
    \end{enumerate}
\end{exercise}

\begin{remark}
    The proof in \textbf{Ex.1.5(ii)} is identical to the one in \textbf{Ex.1.2(ii)}
\end{remark}

\begin{exercise}
    \textbf{(UNFINISHED)}

    By \emph{Proposition 1.11}, if the nilradical is equal to the Jacobson radical,
    since any prime ideal would contain the Jacobson radical, i.e. intersection of all maximal ideals,
    it would contain some of those maixmal ideals.

    Therefore the problem is equivalent to prove that every prime ideals in $A$ is maximal,
    % which is to say the Krull dimension of $A$ is $0$.

    % By \emph{Algebraic Geometry - Hartshorne[Proposition 1.7]}
\end{exercise}

\begin{exercise}
    For any prime ideal $\mathfrak{p}\subseteq A$, consider the map
    $$
    A\longrightarrow A / \mathfrak{p}
    $$

    From the problem we know $\forall x\in A,\exists n>0,\mathrm{s.t.}\ x^{n}=x$.
    So $x(x^{n-1}-1)=0$, and we obtain
    $$
    \overline{x}(\overline{x}^{n-1}-\overline{1})=0
    $$

    We may assume $x\not\in \mathfrak{p}$, i.e. $\overline{x}\neq 0$.
    Since $\mathfrak{p}$ is prime, $A / \mathfrak{p}$ is a integral domain.
    So $\overline{x}^{n-1}=\overline{1}$, $\overline{x}$ has its inverse in $A / \mathfrak{p}$.
    Because $x$ was an arbitary element in $A \backslash \mathfrak{p}$, $A / \mathfrak{p}$ is a field.
    Hence $\mathfrak{p}$ is maximal.
    \qed
\end{exercise}

\begin{exercise}
    Consider the intersection of all prime ideals of $A$, which is the nilradical and obviously minimal with respect to inclusion.
    \qed
\end{exercise}

\begin{exercise}
    Recall that $\rad(\mathfrak{a})$ is the intersection of all prime ideals containing $\mathfrak{a}$.
    The proof can be obtained directly from it.
    \qed
\end{exercise}

\begin{exercise}
    \leavevmode \vspace{-\baselineskip}
    \begin{enumerate}[label=({\alph*})]
        \item  $(i)\Rightarrow(ii):$
            If $A$ has exactly one prime ideal $\mathfrak{p}$, such ideal will be the unique maximal ideal of $A$.

            Recall that, by \emph{Corollary 1.5}, every non-unit of $A$ lies in $\mathfrak{p}$, therefore nilpotent.
            On the other hand, others will be the uni of $A$.
        \item $(ii)\Rightarrow(iii):$ 
            $\forall \text{non-zero }\overline{a}\in \overline{A},\exists a^{-1}\in A\backslash \mathfrak{R},\mathrm{s.t.}\ aa^{-1}=1$.
            And we have $\overline{a}\overline{a^{-1}}=\overline{1}$.
        \item $(iii)\Rightarrow(i):$
            $A / \mathfrak{R}$ is a field if and only if $\mathfrak{R}$ is a maximal ideal.
            Because a nilradial is the intersection of all prime ideals, we can see $\mathfrak{p}$ is unique and maximal in $A$.
    \end{enumerate}
    \qed
\end{exercise}

\begin{exercise}
    \leavevmode \vspace{-\baselineskip}
    \begin{enumerate}[label=({\roman*})]
        \item $2x=4x^2=4x$, so $2x=0$.
        \item For any prime ideal $\mathfrak{p}$, consider the map
            $$
            \varphi:A\longrightarrow A / \mathfrak{p}
            $$
            $$
            x\longrightarrow \overline{x}
            $$
            Since $x^2=x$, $\overline{x}^2=\overline{x},\overline{x}(\overline{x}-1)=0$.
            $A / \mathfrak{p}$ is a integral domain, so $\overline{x}=0$ or $\overline{x}=1$.
            Hence $A / \mathfrak{p}$ only contains two elements, i.e. $\overline{1}$ and $\overline{0}$, and it's a field for sure.
        \item For $I=(x,y)$, consider $z=x+y+xy$, we'll prove $(z)=(x,y)$.

            By definition, $(z)\subseteq (x,y)$.
            Also, notice
            $$
            zx=(x+y+xy)x=x+xy+xy=x
            $$
            So $x\in (z)$.
            In the same way we have $y\in (z)$,
            thus $(z)=(x,y)$.
            Rest of the proof can be done by induction.
    \end{enumerate}
    \qed
\end{exercise}

\begin{exercise}
    
\end{exercise}

\begin{exercise}
    
\end{exercise}

\end{document}
